% SPDX-License-Identifier: CC-BY-4.0
% Copyright (c) 2026 Martin Schröder <info@swedishembedded.com>
\section{A second task family: editing files with Vim, taught from its user manual}\label{sec:vim}

The second family replaces the interpreter-labelled programs with agentic file editing: a request names an outcome, the agent
has one tool that runs Ex commands in a fresh Vim on the task's files and returns Vim's output, and a hidden check grades the
files it leaves. The question of this section is the first step of the loop in that family: do trajectories on tasks that
Splinter authors from a source document alone (Vim~9.1's user manual) teach Qwen3-4B skills that it applies to exams it never
saw, and which property of the training material carries the effect. One generation has been run and examined; it was
registered on 2026-11-26 (splinter \texttt{c319ff4c3}, arm~G added in \texttt{66277f180}) before any arm trained, in
\texttt{.agents/roadmap/source-to-transfer.md}. Everything below is generation~1 unless stated.

\subsection{Pipeline}\label{sec:vim-pipeline}

\begin{figure}[tbp]
\centering
\resizebox{\linewidth}{!}{%
\begin{tikzpicture}[fig, node distance=5mm and 7mm]
  \node[data] (man) {Vim 9.1 user manual\\{\scriptsize 263 items}};
  \node[proc, right=of man, yshift=10mm] (fw) {forward authoring\\{\scriptsize 3 variants per item}};
  \node[proc, right=of fw] (adm) {Vim runs each reference; review, fix};
  \node[data, right=of adm] (ft) {1030 tasks\\{\scriptsize oracle solves 1026}};
  \node[proc, right=of man, yshift=-10mm] (bw) {backward authoring\\{\scriptsize start and end files}};
  \node[proc, right=of bw] (rep) {teacher reaches the end state};
  \node[data, right=of rep] (bt) {467 tasks\\{\scriptsize request written last}};
  \node[proc, right=of $(ft)!0.5!(bt)$, xshift=12mm] (tr) {attempts\\{\scriptsize student 8, teacher 4}};
  \node[data] (ds) at ($(tr |- bt)+(0,-17mm)$) {arm datasets\\{\scriptsize B C D F G}};
  \node[proc, left=12mm of ds] (lora) {LoRA from the base\\{\scriptsize 3 seeds per arm}};
  \node[data, left=12mm of lora] (ex) {100 sealed exams\\{\scriptsize 4 attempts, paired base}};
  \draw[flow] (man.east) -- ++(3mm,0) |- (fw); \draw[flow] (man.east) -- ++(3mm,0) |- (bw);
  \draw[flow] (fw) -- (adm); \draw[flow] (adm) -- (ft);
  \draw[flow] (bw) -- (rep); \draw[flow] (rep) -- (bt);
  \draw[flow] (ft) -| (tr); \draw[flow] (bt) -| (tr);
  \draw[flow] (tr) -- node[edge label, right] {verified, T1 held out} (ds);
  \draw[flow] (ds) -- (lora); \draw[measure] (lora) -- (ex);
\end{tikzpicture}}
\caption{The pipeline as run for generation~1. Rounded boxes run; boxes with a bar are stored. The author and teacher is
Qwen3.8-27B-FP8, the student Qwen3-4B. 10\% of each family's tasks (by task id) are held out of every arm's training (T1). The
exams are hand-written and were never read by any stage.}\label{fig:vim-pipeline}
\end{figure}

\paragraph{Source.} The text of Vim~9.1's user manual (\texttt{usr\_*.txt} as the task image's Vim ships it), with chapter and
section headings added by bench's \texttt{environments/vim/manual\_headings.py}: 222 sections split into 263 items of at most
about 1500 words.

\paragraph{Forward authoring} (bench \texttt{recipes/vim-manual/transfer.yaml}). For each item the 27B author proposes tasks
of one to six Ex-command steps in three material variants, each with a request, starting files and a reference solution.
Vim runs every reference; a reference Vim refuses is repaired once by the author; a reviewer (the same model) checks each
result against its request, refused results are fixed and reviewed again; admission rejects malformed tasks, references that
change nothing or err in Vim, duplicates, and requests that quote the reference's commands. Of 6705 candidate tasks, 1030 were
admitted (682 with one step, 252 with two, 71 with three, 25 with four). The most frequent refusals: a step that changes
nothing (1822), a review that found the result does not match the request (1209), Vim's \texttt{E37} (500), a duplicate
(454). The oracle (the reference replayed by the harness) solved 1026 of the 1030; the four failures were an admission
defect, not task defects (\cref{sec:vim-failed}).

\paragraph{Backward, outcome-first authoring} (bench \texttt{recipes/vim-manual/backward.yaml}). For each item the author
proposes the starting files and the desired end files of a job of several changes, without a request. 771 proposals; 66
refused (file names that are not plain names, carriage returns, duplicates); the teacher then had to reach each of the 705
end states with the Vim tool in 4 attempts, and reached 562. Of those, 523 replay to the end state when Vim reruns the
teacher's commands (30 did not; 9 retyped the files instead of editing them). Only then is the request written, from the
before and after files, and the teacher's verified solution is the reference; 56 requests that quote most of the target or the
solution were refused, leaving 467 tasks, all solved by the oracle.

\paragraph{Trajectories.} Every admitted task was attempted by the student (8 attempts, temperature 0.7, top-p 0.8, top-k 20)
and the teacher (4 attempts) through sven, thinking off, with sven's narrate-before-each-call prompt and a budget of 60 turns
and 16384 completion tokens. \Cref{tab:vim-attempts} gives the pass rates; on the exams the teacher solved 94 of 100 within 4
attempts (pass@1 0.63) and the base student passed none of 1012 counted attempts (calibration, selects nothing).

\begin{table}[tbp]
\centering\small
\begin{tabular}{llrrrr}
\toprule
Material & Model & Counted attempts & Pass rate & Tasks & Tasks solved \\
\midrule
forward & Qwen3-4B & 8204 & 0.100 & 1026 & 231 \\
forward & Qwen3.8-27B-FP8 & 4095 & 0.811 & 1026 & 944 \\
backward & Qwen3-4B & 3688 & 0.031 & 467 & 36 \\
backward & Qwen3.8-27B-FP8 & 1868 & 0.722 & 467 & 410 \\
\bottomrule
\end{tabular}
\caption{Attempts on the authored tasks (\texttt{material/funnel.json} of the evidence). The backward tasks are harder for both
models and closer to the exams in difficulty.}\label{tab:vim-attempts}
\end{table}

\begin{figure}[tbp]
\centering
\begin{tikzpicture}
\begin{groupplot}[group style={group size=2 by 1, horizontal sep=24mm}, paper, width=0.34\linewidth, height=50mm,
    xmode=log, xmin=10, xmax=40000, ytick=data, y tick label style={font=\scriptsize},
    every axis plot/.append style={xbar, /pgf/bar width=2.4mm},
    xmajorgrids, ymajorgrids=false, nodes near coords, nodes near coords style={font=\scriptsize},
    point meta=explicit symbolic, enlarge y limits=0.08, table/col sep=comma]
\nextgroupplot[title={forward}, yticklabels from table={generated/vim-transfer-funnel-forward.csv}{step}]
  \addplot[draw=black!70, fill=okBlue!70] table[col sep=comma, x=count, y=y, meta=count] {generated/vim-transfer-funnel-forward.csv};
\nextgroupplot[title={backward (outcome first)}, yticklabels from table={generated/vim-transfer-funnel-backward.csv}{step}]
  \addplot[draw=black!70, fill=okOrange!80] table[col sep=comma, x=count, y=y, meta=count] {generated/vim-transfer-funnel-backward.csv};
\end{groupplot}
\end{tikzpicture}
\caption{Counts through authoring (log scale; \texttt{material/funnel.json} of the evidence). ``Proposal sets'' are the
author's outputs per item and variant; ``teacher/student solved'' count tasks with at least one verified attempt.}
\label{fig:vim-funnel}
\end{figure}

\subsection{Registered design}\label{sec:vim-registration}

\paragraph{Arms.} All arms start from the same base Qwen3-4B and train a LoRA adapter (rank 16, alpha 32, learning rate
$2\times10^{-4}$, 800 steps of 4 records, block 16384; the checkpoint is chosen by held-out validation loss every 40 steps
with patience 4), seeds 1, 2 and 3. A is the base. B trains on every verified successful attempt of teacher and student on the
forward tasks, every action supervised. C keeps at most one attempt per task, the one with the fewest actions. D is B with an
action supervised only if at least one of its tool calls succeeded (an all-error action stays in context, unsupervised).
F, the negative control, is B with each record's request replaced by another task's (a seeded derangement): the same records
and tool-use format with no request-to-action mapping. G applies B's selection to the backward tasks; it was added before any
arm trained, as a dated deviation.

\paragraph{Measurements.} T2/T4: the 100 held-out course exams of bench's \texttt{tasks/vim} (hand-written, unseen material,
compositions of several ideas), 4 attempts per exam per arm-seed with the practice sampling and budget, each arm-seed against a
base run in the same evaluation stage. T1: the held-out authored tasks, 4 attempts each. An attempt the infrastructure ended
is never a failure; a stage with more than 2\% such attempts is repeated after its cause is fixed.

\paragraph{Decision rule (as registered).} Unit: an exam. For arms X and Y, $d_t$ is X's passes over its 3 seeds $\times$ 4
attempts minus Y's over 12 attempts (for A, the base runs paired with X); one-sided exact sign test over exams with
$d_t\neq0$. Primary family, Holm at $\alpha=0.05$: B, C, D, G each greater than A on the exams. The claim that trajectories
teach transferable skills holds when at least one of B, C, D beats A in the primary family and the same arm beats F
(one-sided sign test, $\alpha=0.05$). Secondary, with task-bootstrap 95\% intervals and never claims when not significant:
G, D and C each against B, every arm on T1, per-seed rates. No exam result selects a checkpoint, recipe, seed or data mix.

\subsection{Deviations}\label{sec:vim-deviations}

\begin{table}[h]
\centering\footnotesize
\begin{tabular}{p{17mm}p{64mm}p{60mm}}
\toprule
Date & Change & Why \\
\midrule
2026-11-26 & arm G added; the primary family grows to four comparisons & outcome-first authoring was proposed while the material
run was in progress; added before any arm trained \\
2026-11-26 & source is the manual's text (\texttt{usr\_*.txt}), not an OCR of its PDF & the OCR copy was lost with a node \\
2026-11-26 & the material run was forked once; its authoring, review, task and oracle stages reused, all student and teacher
attempts collected again & a node ran out of memory and the serving fleet hung; no attempt of the first run is used \\
2026-11-26 & servers run without speculative decoding & 7 of 16 servers hung on a device kernel with it; the sampled
distribution is unchanged \\
2026-11-26 & each arm-seed trains on one card and evaluates through its node's four cards & a router does not route training \\
2026-11-27 & the operational definition (\cref{sec:introduction}) written after generation~1 was examined, before any verdict
against it & requested so that the claim is scoped by a fixed definition; adds no registered comparison \\
\bottomrule
\end{tabular}
\caption{Deviations from the registration, as dated in the ledger.}\label{tab:vim-deviations}
\end{table}

\subsection{Results of generation 1}\label{sec:vim-results}

\paragraph{Run.} All 15 arm-seeds trained and were examined. Infrastructure ended at most 11 of an evaluation stage's 800
attempts (1.4\%), all on the base side; the bound held in every stage. B, D and F ran all 800 steps; C stopped early and kept
step 320 or 360, G kept step 520 in every seed. The optimizer budget is therefore not equal across arms: B, D and F saw 3200
records (0.95 passes over 3359), G 2080 (1.9 passes over 1100), C 1280--1440 (1.7--1.9 passes over 771). Held-out validation
loss fell from 2.35--2.81 (base) to 0.57--0.68; seeds of one arm differ by less than 0.01.

\begin{table}[tbp]
\centering\footnotesize
\setlength{\tabcolsep}{3.5pt}
\resizebox{\linewidth}{!}{%
\begin{tabular}{lrrlrrll}
\toprule
Arm & Train records & Teacher / student & Passes of 1200 (by seed) & Exams solved & Base & vs A: W/L, $p$ & vs F: W/L, $p$ \\
\midrule
\inputrows{generated/vim-transfer-table.tex}
\bottomrule
\end{tabular}}
\caption{Generation~1 on the 100 exams (registered analysis; evidence file \nolinkurl{decision/decision-eval-exam.json}). ``Base'' is the
passes of the paired base runs over their counted attempts. One-sided exact sign tests over exams; all four primary
comparisons hold under Holm (adjusted thresholds 0.0125 to 0.05). F's comparison with A is not in the registered family;
F and its paired base both pass none.}\label{tab:vim-results}
\end{table}

\begin{figure}[tbp]
\centering
\begin{tikzpicture}
\begin{axis}[paper, width=0.72\linewidth, height=55mm, xmin=0.4, xmax=5.6, ymin=0, ymax=0.075,
    xtick={1,2,3,4,5}, xticklabels={B,C,D,F,G}, ylabel={exam pass rate}, yticklabel style={/pgf/number format/fixed},
    scaled y ticks=false, legend pos=north west]
  \addplot[only marks, mark=*, mark size=1.5pt, okBlue, error bars/.cd, y dir=both, y explicit]
    table[col sep=comma, x=x, y=rate, y error plus expr=\thisrow{hi}-\thisrow{rate}, y error minus expr=\thisrow{rate}-\thisrow{lo}]
    {generated/vim-transfer-seeds.csv};
  \addlegendentry{seed (Wilson 95\%)}
  \addplot[only marks, mark=diamond*, mark size=3pt, okVermillion, error bars/.cd, y dir=both, y explicit]
    table[col sep=comma, x expr=\thisrow{x}+0.36, y=rate, y error plus expr=\thisrow{hi}-\thisrow{diff},
      y error minus expr=\thisrow{diff}-\thisrow{lo}] {generated/vim-transfer-arms.csv};
  \addlegendentry{arm, 3 seeds (task bootstrap)}
  \addplot[okGrey, densely dotted, thick, domain=0.4:5.6] {0};
\end{axis}
\end{tikzpicture}
\caption{Exam pass rate per arm and seed (400 attempts each) with Wilson intervals, which treat attempts as independent and
are therefore too narrow for exams that cluster; and the pooled arm (1200 attempts) with the task-bootstrap interval of its
difference from the paired base. Every base run passed none (dotted).}\label{fig:vim-seeds}
\end{figure}

\paragraph{Registered verdict.} Every primary comparison holds: B (10 exams won, 0 lost, $p=0.00098$), C (6/0, $p=0.016$),
D (11/0, $p=0.00049$) and G (24/0, $p=6.0\times10^{-8}$) each beat the base under Holm, and each beats the shuffled-request
control F (B 10/0, C 6/0, D 11/0, G 25/0). The registered claim holds for generation~1: trajectories on tasks authored from
the manual alone teach the 4B student skills it applies to exams it never saw, and F, trained on the same records and format
with deranged requests, stays at the base's zero, so the gain comes from the pairing of request and actions, not from exposure
to the tool's format. Every seed of B, C, D and G passes at least two exam attempts; every seed of F passes none
(\cref{fig:vim-seeds}).

\paragraph{Effect size.} The effects are small in absolute terms. G passes 42 of 1200 exam attempts (3.5\%; difference from
its paired base $+0.035$, 95\% interval $[+0.021, +0.051]$) and solves 25 of the 100 exams at least once; B and D pass 1.8\%
($+0.019$ $[+0.008, +0.032]$), C 0.8\% ($+0.008$ $[+0.003, +0.017]$). The teacher passes 63\% of first attempts on the same
exams.

\paragraph{Secondary comparisons} (two-sided, not claims of the registered family):

\begin{center}\footnotesize
\begin{tabular}{lrll}
\toprule
Comparison & Wins/losses & $p$ & Pass-rate difference [95\% interval] \\
\midrule
\inputrows{generated/vim-transfer-secondary.tex}
\bottomrule
\end{tabular}
\end{center}

G beats B although it trained on a third of B's records (1100 against 3359) from 348 tasks against 857: on this exam,
outcome-first tasks transfer more per record than section-by-section tasks. The comparison is not controlled for the
optimizer budget per record (G saw each record about twice, B about once) or for difficulty (the backward tasks are harder,
\cref{tab:vim-attempts}), so it does not say which of these causes it. Masking all-error actions (D) changed nothing
measurable. Keeping one solution per task (C) taught less than keeping all verified solutions; C also trained on fewer
records, so this does not separate diversity from volume.

\paragraph{Teacher and student records.} Of B's 3359 training records, 2699 come from the teacher and 660 from the student
(G: 1014 and 86; C: 629 and 142), by each record's \texttt{model} field in the dataset's lineage. The arms therefore do not
separate what the teacher's trajectories teach from what the student's own verified attempts teach. A student-only arm S
(B's rule restricted to the 4B's own verified successes, 3 seeds) was started after these results to separate them; its
result is open.

\paragraph{T1, held-out authored tasks} (secondary). Arm C on the 94 held-out forward tasks: 642 of 1128 attempts (0.569)
against its paired base's 113 of 1126 (0.100); 79 tasks won, 3 lost (one-sided $p=1.9\times10^{-20}$), difference $+0.470$
$[+0.397, +0.543]$. In-distribution, the gain is an order of magnitude larger than on the exams. T1 for B, D, F and G was still
being measured when this was written and is open. G's T1 tasks are the held-out backward tasks, a different set from the
others'.

\subsection{Generations, retention, and the operational definition}\label{sec:vim-verdict}

\begin{figure}[tbp]
\centering
\begin{tikzpicture}
\begin{axis}[paper, width=0.62\linewidth, height=52mm, xmin=-0.3, xmax=3.3, ymin=0, ymax=0.06, xtick={0,1,2,3},
    xticklabels={base,gen.~1,gen.~2,gen.~3}, ylabel={exam pass rate}, yticklabel style={/pgf/number format/fixed},
    scaled y ticks=false, legend pos=north west, legend columns=2]
  \fill[black!6] (axis cs:1.5,0) rectangle (axis cs:3.3,0.06);
  \node[font=\scriptsize, text=black!60] at (axis cs:2.4,0.052) {not yet run};
  \addplot[okOrange, thick, mark=triangle*, discard if not={arm}{G}, error bars/.cd, y dir=both, y explicit]
    table[col sep=comma, x=generation, y=rate, y error plus expr=\thisrow{hi}-\thisrow{rate},
      y error minus expr=\thisrow{rate}-\thisrow{lo}] {generated/vim-transfer-generations.csv};
  \addlegendentry{G}
  \addplot[okBlue, thick, mark=*, discard if not={arm}{B}, error bars/.cd, y dir=both, y explicit]
    table[col sep=comma, x=generation, y=rate, y error plus expr=\thisrow{hi}-\thisrow{rate},
      y error minus expr=\thisrow{rate}-\thisrow{lo}] {generated/vim-transfer-generations.csv};
  \addlegendentry{B}
  \addplot[okGreen, thick, dashed, mark=square*, discard if not={arm}{D}]
    table[col sep=comma, x=generation, y=rate] {generated/vim-transfer-generations.csv};
  \addlegendentry{D}
  \addplot[okGrey, thick, densely dotted, mark=diamond*, discard if not={arm}{F}]
    table[col sep=comma, x=generation, y=rate] {generated/vim-transfer-generations.csv};
  \addlegendentry{F (control)}
\end{axis}
\end{tikzpicture}
\caption{Exam pass rate by generation (Wilson 95\% intervals over attempts). One generation exists; later generations are
open.}\label{fig:vim-generations}
\end{figure}

\paragraph{Retention.} No forgetting gate was run: the arms were measured on the Vim exams and the authored hold-out only,
not on tasks of earlier skills (for instance the Lua family of \cref{sec:results} or general instruction following). Retention
is unmeasured.

\paragraph{Verdict against the operational definition} (\cref{sec:introduction}). (a) Held-out gain: holds for generation~1
against the base (every primary comparison); no later generation exists, so the requirement of three successive generations
is unmet. (b) No forgetting: unmet, not measured. (c) Robust and content-borne: holds for generation~1 (all three seeds of
every arm but F improve; every arm beats F). (d) Restartable: unmet; the material run was forked by hand after a node ran out
of memory and the serving fleet hung, and several pipeline defects were fixed during the run (\cref{sec:vim-failed}). The
supported claim is therefore one step of the loop in a second task family, not continuous self-improvement. In the Lua
family, RSI4 (the companion paper's S4~\cite{frontier2026}) held part~(a) over three generations in four of four lineages, with
no forgetting gate and no shuffled control.

\subsection{What failed or was fixed during the run}\label{sec:vim-failed}

\begin{itemize}
\item Four of 1030 admitted tasks failed the oracle: admission read Vim's output in text mode (turning CR~LF into LF) and ran
  Vim under a locale the task image lacks (so the euro sign was written as a different byte). Fixed in bench \texttt{f73ae591}
  (knowledge 086) after the material was made; the four tasks were never attempted (attempts require an oracle solution).
\item Sandbox images held in memory exhausted a node (splinter \texttt{851216e17}, knowledge 225), forcing the fork of the
  material run; a resumed batch could not reconcile (\texttt{5fb824d6b}); attempts lost to interruptions spent retries
  (\texttt{cc5581333}).
\item On-disk sandbox copies took 436~s to start (shared lower tree plus overlay, \texttt{9841886df}, knowledge 227); the
  gateway's connection pool, Harbor's job lock and a per-command lookup took about 50 of 70~s per trial (\texttt{9415eb39b},
  \texttt{2b62ff444}, \texttt{de1149b0a}, \texttt{418046bd7}, knowledge 228).
\item Brain's router stranded requests behind hung upstreams (\texttt{ae72fc212}) and packed one model on few cards
  (\texttt{b726c045b}); Qwen3 serving plans withdrew a model instead of fitting free memory (\texttt{f74314255}); a CUDA lane
  can hang while \texttt{/health} reports ok (brain knowledge 602; no watchdog yet).
\item Of the 705 backward jobs the teacher attempted, 1498 of 2820 attempts ran out of wall-clock time; 143 end states were
  never reached and are not tasks.
\end{itemize}

\subsection{Evidence and reproduction}\label{sec:vim-evidence}
\begin{sloppypar}

Aggregate evidence is in the private Hugging Face dataset \nolinkurl{swedishembedded/source-to-transfer-evidence} (arm datasets
with lineage, per arm-seed recipes, locks and frozen revisions, training reports, exam aggregates, the decision and secondary
outputs, the authoring funnel, the material runs' manifests, and the analysis code); the 15 adapters are in
\nolinkurl{swedishembedded/vim-manual-transfer-adapters}; the authored tasks and trajectories in
\nolinkurl{swedishembedded/vim-headless-mastery-manual} and \nolinkurl{swedishembedded/vim-headless-mastery-sft-manual}. Per-exam
outcomes are item-level exam data and stay with the examiner's private evidence. Frozen builds: splinter \texttt{f6661e34},
sven \texttt{704496d6}, brain \texttt{b2f46a22}; recipes from bench \texttt{fac113db} (forward), \texttt{83e1b4f8}
(backward) and \texttt{662ac1d6} (arms). The decision is recomputed by \texttt{decide.py RESULTS eval-exam} and
\texttt{secondary.py RESULTS eval-exam} from the per-task outcomes \texttt{export-results.py} writes; the tables and figures
of this section by \texttt{scripts/vim\_transfer\_data.py} from the evidence files copied into \texttt{data/vim-transfer}
(\texttt{make vim-data EVIDENCE=DIR}).
\end{sloppypar}
